\chapter{30 câu nói nhân sinh cực sâu}
\markboth{30 câu nói nhân sinh cực sâu}{30 Deep Life Sayings}


\section{Sống là chấp nhận những gì không như ý}
\begin{truyen}{Sống là chấp nhận những gì không như ý}{Life lesson}
\chuhoa{N}{gười khôn không đòi đời phải đúng ý mình. Sống là chấp nhận những gì không như ý.}
\textit{(The wise don't demand life to go their way. To live is to accept what isn't as wished.)}
\end{truyen}

\begin{baihoc}
Chấp nhận không phải đầu hàng.
\textit{(To accept is not to surrender.)}
\end{baihoc}

\begin{ghinhoanh}
\textbf{Short English to remember:}

To accept is not to surrender.

\end{ghinhoanh}

\begin{tuhocnhanh}
1. Câu chuyện này gợi cho bạn suy nghĩ gì? \textit{What does this story make you think?}

2. Bạn sẽ nhớ bài học nào nhất? \textit{Which lesson will you remember most?}
\end{tuhocnhanh}
\ngancach


\section{Đau đớn dạy ta điều hạnh phúc không dạy được}
\begin{truyen}{Đau đớn dạy ta điều hạnh phúc không dạy được}{Life lesson}
\chuhoa{N}{hững bài học lớn thường đến từ đau. Đau đớn dạy ta điều hạnh phúc không dạy được.}
\textit{(The biggest lessons often come from pain. Pain teaches what happiness cannot.)}
\end{truyen}

\begin{baihoc}
Đau là thầy.
\textit{(Pain is a teacher.)}
\end{baihoc}

\begin{ghinhoanh}
\textbf{Short English to remember:}

Pain is a teacher.

\end{ghinhoanh}

\begin{tuhocnhanh}
1. Câu chuyện này gợi cho bạn suy nghĩ gì? \textit{What does this story make you think?}

2. Bạn sẽ nhớ bài học nào nhất? \textit{Which lesson will you remember most?}
\end{tuhocnhanh}
\ngancach


\section{Im lặng đôi khi là câu trả lời cao nhất}
\begin{truyen}{Im lặng đôi khi là câu trả lời cao nhất}{Life lesson}
\chuhoa{K}{hông phải lúc nào cũng cần lời. Im lặng đôi khi là câu trả lời cao nhất.}
\textit{(Not every moment needs words. Silence is sometimes the highest answer.)}
\end{truyen}

\begin{baihoc}
Im lặng có sức nặng.
\textit{(Silence has weight.)}
\end{baihoc}

\begin{ghinhoanh}
\textbf{Short English to remember:}

Silence has weight.

\end{ghinhoanh}

\begin{tuhocnhanh}
1. Câu chuyện này gợi cho bạn suy nghĩ gì? \textit{What does this story make you think?}

2. Bạn sẽ nhớ bài học nào nhất? \textit{Which lesson will you remember most?}
\end{tuhocnhanh}
\ngancach


\section{Người hiểu mình nhất thường là người đã làm mình đau}
\begin{truyen}{Người hiểu mình nhất thường là người đã làm mình đ...}{Life lesson}
\chuhoa{K}{ẻ từng làm ta đau đôi khi là người hiểu ta nhất. Người hiểu mình nhất thường là người đã làm mình đau.}
\textit{(The one who hurt us sometimes understands us best. Those who understand you most have often hurt you.)}
\end{truyen}

\begin{baihoc}
Đau và hiểu đôi khi đi cùng.
\textit{(Pain and understanding sometimes go together.)}
\end{baihoc}

\begin{ghinhoanh}
\textbf{Short English to remember:}

Pain and understanding sometimes go together.

\end{ghinhoanh}

\begin{tuhocnhanh}
1. Câu chuyện này gợi cho bạn suy nghĩ gì? \textit{What does this story make you think?}

2. Bạn sẽ nhớ bài học nào nhất? \textit{Which lesson will you remember most?}
\end{tuhocnhanh}
\ngancach


\section{Thời gian không chữa lành — ta học sống chung với vết sẹo}
\begin{truyen}{Thời gian không chữa lành — ta học sống chung với ...}{Life lesson}
\chuhoa{V}{ết thương có thể không biến mất. Thời gian không chữa lành — ta học sống chung với vết sẹo.}
\textit{(Wounds may not disappear. Time doesn't heal — we learn to live with the scar.)}
\end{truyen}

\begin{baihoc}
Sống chung với quá khứ.
\textit{(Living with the past.)}
\end{baihoc}

\begin{ghinhoanh}
\textbf{Short English to remember:}

Living with the past.

\end{ghinhoanh}

\begin{tuhocnhanh}
1. Câu chuyện này gợi cho bạn suy nghĩ gì? \textit{What does this story make you think?}

2. Bạn sẽ nhớ bài học nào nhất? \textit{Which lesson will you remember most?}
\end{tuhocnhanh}
\ngancach


\section{Cho đi không cần nhận lại — nhưng nhận lại được là phúc}
\begin{truyen}{Cho đi không cần nhận lại — nhưng nhận lại được là...}{Life lesson}
\chuhoa{A}{nh cho mà không đòi. Nhưng khi được nhận lại, anh biết đó là phúc. Cho đi không cần nhận lại — nhưng nhận lại được là phúc.}
\textit{(He gave without demanding. But when he received back, he knew it was a blessing.)}
\end{truyen}

\begin{baihoc}
Cho là tự nguyện.
\textit{(Giving is voluntary.)}
\end{baihoc}

\begin{ghinhoanh}
\textbf{Short English to remember:}

Giving is voluntary.

\end{ghinhoanh}

\begin{tuhocnhanh}
1. Câu chuyện này gợi cho bạn suy nghĩ gì? \textit{What does this story make you think?}

2. Bạn sẽ nhớ bài học nào nhất? \textit{Which lesson will you remember most?}
\end{tuhocnhanh}
\ngancach


\section{Khiêm tốn là khi mình giỏi mà không cần chứng minh}
\begin{truyen}{Khiêm tốn là khi mình giỏi mà không cần chứng minh}{Life lesson}
\chuhoa{N}{gười thật sự khiêm tốn không cần khoe. Khiêm tốn là khi mình giỏi mà không cần chứng minh.}
\textit{(The truly humble don't need to show off. Humility is being good without having to prove it.)}
\end{truyen}

\begin{baihoc}
Khiêm tốn là sức mạnh.
\textit{(Humility is strength.)}
\end{baihoc}

\begin{ghinhoanh}
\textbf{Short English to remember:}

Humility is strength.

\end{ghinhoanh}

\begin{tuhocnhanh}
1. Câu chuyện này gợi cho bạn suy nghĩ gì? \textit{What does this story make you think?}

2. Bạn sẽ nhớ bài học nào nhất? \textit{Which lesson will you remember most?}
\end{tuhocnhanh}
\ngancach


\section{Tha thứ là món quà cho chính mình}
\begin{truyen}{Tha thứ là món quà cho chính mình}{Life lesson}
\chuhoa{A}{nh tha thứ không phải vì người kia xứng đáng. Anh tha để mình nhẹ. Tha thứ là món quà cho chính mình.}
\textit{(He forgave not because the other deserved it. He forgave so he could be light. Forgiveness is a gift to yourself.)}
\end{truyen}

\begin{baihoc}
Tha thứ để bước tiếp.
\textit{(Forgive to move on.)}
\end{baihoc}

\begin{ghinhoanh}
\textbf{Short English to remember:}

Forgive to move on.

\end{ghinhoanh}

\begin{tuhocnhanh}
1. Câu chuyện này gợi cho bạn suy nghĩ gì? \textit{What does this story make you think?}

2. Bạn sẽ nhớ bài học nào nhất? \textit{Which lesson will you remember most?}
\end{tuhocnhanh}
\ngancach


\section{Không ai hiểu hết nỗi đau của người khác}
\begin{truyen}{Không ai hiểu hết nỗi đau của người khác}{Life lesson}
\chuhoa{Đ}{ừng nói "tôi hiểu" khi chưa đi qua. Không ai hiểu hết nỗi đau của người khác.}
\textit{(Don't say "I understand" until you've been there. No one fully understands another's pain.)}
\end{truyen}

\begin{baihoc}
Lắng nghe thay vì phán xét.
\textit{(Listen instead of judge.)}
\end{baihoc}

\begin{ghinhoanh}
\textbf{Short English to remember:}

Listen instead of judge.

\end{ghinhoanh}

\begin{tuhocnhanh}
1. Câu chuyện này gợi cho bạn suy nghĩ gì? \textit{What does this story make you think?}

2. Bạn sẽ nhớ bài học nào nhất? \textit{Which lesson will you remember most?}
\end{tuhocnhanh}
\ngancach


\section{Sống chậm lại mới thấy mình đang sống}
\begin{truyen}{Sống chậm lại mới thấy mình đang sống}{Life lesson}
\chuhoa{C}{hạy mãi, anh quên mình đang chạy vì đâu. Sống chậm lại mới thấy mình đang sống.}
\textit{(Running forever, he forgot why he was running. Slow down to see that you're alive.)}
\end{truyen}

\begin{baihoc}
Chậm là tỉnh.
\textit{(Slow is awake.)}
\end{baihoc}

\begin{ghinhoanh}
\textbf{Short English to remember:}

Slow is awake.

\end{ghinhoanh}

\begin{tuhocnhanh}
1. Câu chuyện này gợi cho bạn suy nghĩ gì? \textit{What does this story make you think?}

2. Bạn sẽ nhớ bài học nào nhất? \textit{Which lesson will you remember most?}
\end{tuhocnhanh}
\ngancach


\section{Người mạnh không phải không khóc — mà khóc xong vẫn đứng dậy}
\begin{truyen}{Người mạnh không phải không khóc — mà khóc xong vẫ...}{Life lesson}
\chuhoa{K}{hóc không làm anh yếu. Bỏ cuộc mới làm. Người mạnh không phải không khóc — mà khóc xong vẫn đứng dậy.}
\textit{(Crying doesn't make you weak. Giving up does. The strong aren't those who don't cry — they're those who still stand after.)}
\end{truyen}

\begin{baihoc}
Khóc rồi đứng dậy.
\textit{(Cry then stand.)}
\end{baihoc}

\begin{ghinhoanh}
\textbf{Short English to remember:}

Cry then stand.

\end{ghinhoanh}

\begin{tuhocnhanh}
1. Câu chuyện này gợi cho bạn suy nghĩ gì? \textit{What does this story make you think?}

2. Bạn sẽ nhớ bài học nào nhất? \textit{Which lesson will you remember most?}
\end{tuhocnhanh}
\ngancach


\section{Tương lai không chờ ai — nhưng quá khứ có thể níu chân}
\begin{truyen}{Tương lai không chờ ai — nhưng quá khứ có thể níu ...}{Life lesson}
\chuhoa{A}{nh không thể đổi quá khứ. Anh chỉ có thể đổi cách mình đi tiếp. Tương lai không chờ ai — nhưng quá khứ có thể níu chân.}
\textit{(You can't change the past. You can only change how you go on. The future waits for no one — but the past can hold you back.)}
\end{truyen}

\begin{baihoc}
Buông quá khứ để bước.
\textit{(Let go of the past to step.)}
\end{baihoc}

\begin{ghinhoanh}
\textbf{Short English to remember:}

Let go of the past to step.

\end{ghinhoanh}

\begin{tuhocnhanh}
1. Câu chuyện này gợi cho bạn suy nghĩ gì? \textit{What does this story make you think?}

2. Bạn sẽ nhớ bài học nào nhất? \textit{Which lesson will you remember most?}
\end{tuhocnhanh}
\ngancach


\section{Tử tế với người không tử tế với mình — đó mới là tu}
\begin{truyen}{Tử tế với người không tử tế với mình — đó mới là t...}{Life lesson}
\chuhoa{D}{ễ tử tế với người tử tế. Khó là tử tế với kẻ không tử tế với mình. Tử tế với người không tử tế với mình — đó mới là tu.}
\textit{(It's easy to be kind to the kind. The hard part is being kind to those who aren't. Kindness to the unkind — that's the practice.)}
\end{truyen}

\begin{baihoc}
Tu là trong khó.
\textit{(Practice is in the hard.)}
\end{baihoc}

\begin{ghinhoanh}
\textbf{Short English to remember:}

Practice is in the hard.

\end{ghinhoanh}

\begin{tuhocnhanh}
1. Câu chuyện này gợi cho bạn suy nghĩ gì? \textit{What does this story make you think?}

2. Bạn sẽ nhớ bài học nào nhất? \textit{Which lesson will you remember most?}
\end{tuhocnhanh}
\ngancach


\section{Cô đơn dạy ta biết trân trọng khi có người bên cạnh}
\begin{truyen}{Cô đơn dạy ta biết trân trọng khi có người bên cạn...}{Life lesson}
\chuhoa{K}{hi anh cô đơn, anh mới hiểu giá trị của "có ai đó". Cô đơn dạy ta biết trân trọng khi có người bên cạnh.}
\textit{(When you're alone you understand the value of "someone there". Loneliness teaches us to treasure company.)}
\end{truyen}

\begin{baihoc}
Cô đơn là bài học.
\textit{(Loneliness is a lesson.)}
\end{baihoc}

\begin{ghinhoanh}
\textbf{Short English to remember:}

Loneliness is a lesson.

\end{ghinhoanh}

\begin{tuhocnhanh}
1. Câu chuyện này gợi cho bạn suy nghĩ gì? \textit{What does this story make you think?}

2. Bạn sẽ nhớ bài học nào nhất? \textit{Which lesson will you remember most?}
\end{tuhocnhanh}
\ngancach


\section{Thất bại không định nghĩa em — cách em đứng dậy mới định nghĩa}
\begin{truyen}{Thất bại không định nghĩa em — cách em đứng dậy mớ...}{Life lesson}
\chuhoa{N}{gã thì ai cũng ngã. Khác nhau ở chỗ đứng dậy thế nào. Thất bại không định nghĩa em — cách em đứng dậy mới định nghĩa.}
\textit{(Everyone falls. The difference is how you get up. Failure doesn't define you — how you get up does.)}
\end{truyen}

\begin{baihoc}
Đứng dậy là định nghĩa.
\textit{(Getting up is the definition.)}
\end{baihoc}

\begin{ghinhoanh}
\textbf{Short English to remember:}

Getting up is the definition.

\end{ghinhoanh}

\begin{tuhocnhanh}
1. Câu chuyện này gợi cho bạn suy nghĩ gì? \textit{What does this story make you think?}

2. Bạn sẽ nhớ bài học nào nhất? \textit{Which lesson will you remember most?}
\end{tuhocnhanh}
\ngancach


\section{Đừng so sánh bản thân với người khác — so với chính mình hôm qua}
\begin{truyen}{Đừng so sánh bản thân với người khác — so với chín...}{Life lesson}
\chuhoa{S}{o với người khác chỉ thêm mệt. So với mình hôm qua mới thấy tiến. Đừng so sánh bản thân với người khác — so với chính mình hôm qua.}
\textit{(Comparing with others only tires you. Comparing with yesterday's you shows progress. Don't compare yourself to others — compare to yesterday's you.)}
\end{truyen}

\begin{baihoc}
So với mình.
\textit{(Compare with yourself.)}
\end{baihoc}

\begin{ghinhoanh}
\textbf{Short English to remember:}

Compare with yourself.

\end{ghinhoanh}

\begin{tuhocnhanh}
1. Câu chuyện này gợi cho bạn suy nghĩ gì? \textit{What does this story make you think?}

2. Bạn sẽ nhớ bài học nào nhất? \textit{Which lesson will you remember most?}
\end{tuhocnhanh}
\ngancach


\section{Lời nói có thể giết — cũng có thể cứu}
\begin{truyen}{Lời nói có thể giết — cũng có thể cứu}{Life lesson}
\chuhoa{M}{ột câu có thể làm người ta gục. Một câu khác có thể nâng họ dậy. Lời nói có thể giết — cũng có thể cứu.}
\textit{(One sentence can bring someone down. Another can lift them up. Words can kill — and can save.)}
\end{truyen}

\begin{baihoc}
Nói có ý thức.
\textit{(Speak with care.)}
\end{baihoc}

\begin{ghinhoanh}
\textbf{Short English to remember:}

Speak with care.

\end{ghinhoanh}

\begin{tuhocnhanh}
1. Câu chuyện này gợi cho bạn suy nghĩ gì? \textit{What does this story make you think?}

2. Bạn sẽ nhớ bài học nào nhất? \textit{Which lesson will you remember most?}
\end{tuhocnhanh}
\ngancach


\section{Sức khỏe mất rồi mới biết nó đáng giá hơn tiền}
\begin{truyen}{Sức khỏe mất rồi mới biết nó đáng giá hơn tiền}{Life lesson}
\chuhoa{A}{nh đổi sức khỏe lấy tiền. Đến khi bệnh, anh đổi tiền lại sức khỏe — không đủ. Sức khỏe mất rồi mới biết nó đáng giá hơn tiền.}
\textit{(He traded health for money. When he got sick he traded money back for health — not enough. You only know health's value when it's gone.)}
\end{truyen}

\begin{baihoc}
Sức khỏe là vốn.
\textit{(Health is capital.)}
\end{baihoc}

\begin{ghinhoanh}
\textbf{Short English to remember:}

Health is capital.

\end{ghinhoanh}

\begin{tuhocnhanh}
1. Câu chuyện này gợi cho bạn suy nghĩ gì? \textit{What does this story make you think?}

2. Bạn sẽ nhớ bài học nào nhất? \textit{Which lesson will you remember most?}
\end{tuhocnhanh}
\ngancach


\section{Gia đình là nơi về — đừng để khi mất mới về}
\begin{truyen}{Gia đình là nơi về — đừng để khi mất mới về}{Life lesson}
\chuhoa{A}{nh bận, ít về. Một ngày không còn ai để về. Gia đình là nơi về — đừng để khi mất mới về.}
\textit{(He was busy, rarely home. One day there was no one to go home to. Family is where you return — don't wait until it's gone.)}
\end{truyen}

\begin{baihoc}
Về khi còn có thể.
\textit{(Return while you can.)}
\end{baihoc}

\begin{ghinhoanh}
\textbf{Short English to remember:}

Return while you can.

\end{ghinhoanh}

\begin{tuhocnhanh}
1. Câu chuyện này gợi cho bạn suy nghĩ gì? \textit{What does this story make you think?}

2. Bạn sẽ nhớ bài học nào nhất? \textit{Which lesson will you remember most?}
\end{tuhocnhanh}
\ngancach


\section{Học từ sai lầm của người khác — đỡ phải trả giá}
\begin{truyen}{Học từ sai lầm của người khác — đỡ phải trả giá}{Life lesson}
\chuhoa{A}{nh có thể học từ sai lầm của mình. Khôn hơn là học từ sai lầm của người khác. Học từ sai lầm của người khác — đỡ phải trả giá.}
\textit{(You can learn from your own mistakes. Wiser is learning from others'. Learn from others' mistakes — so you pay less.)}
\end{truyen}

\begin{baihoc}
Học từ người khác.
\textit{(Learn from others.)}
\end{baihoc}

\begin{ghinhoanh}
\textbf{Short English to remember:}

Learn from others.

\end{ghinhoanh}

\begin{tuhocnhanh}
1. Câu chuyện này gợi cho bạn suy nghĩ gì? \textit{What does this story make you think?}

2. Bạn sẽ nhớ bài học nào nhất? \textit{Which lesson will you remember most?}
\end{tuhocnhanh}
\ngancach


\section{Kiên nhẫn không phải chờ thụ động — mà chờ trong hành động}
\begin{truyen}{Kiên nhẫn không phải chờ thụ động — mà chờ trong h...}{Life lesson}
\chuhoa{A}{nh không ngồi chờ. Anh vừa làm vừa chờ. Kiên nhẫn không phải chờ thụ động — mà chờ trong hành động.}
\textit{(He didn't just sit and wait. He acted while waiting. Patience isn't passive waiting — it's waiting in action.)}
\end{truyen}

\begin{baihoc}
Chờ mà vẫn làm.
\textit{(Wait while doing.)}
\end{baihoc}

\begin{ghinhoanh}
\textbf{Short English to remember:}

Wait while doing.

\end{ghinhoanh}

\begin{tuhocnhanh}
1. Câu chuyện này gợi cho bạn suy nghĩ gì? \textit{What does this story make you think?}

2. Bạn sẽ nhớ bài học nào nhất? \textit{Which lesson will you remember most?}
\end{tuhocnhanh}
\ngancach


\section{Biết ơn những gì nhỏ — sẽ thấy mình giàu}
\begin{truyen}{Biết ơn những gì nhỏ — sẽ thấy mình giàu}{Life lesson}
\chuhoa{A}{nh tập biết ơn bữa cơm, người thân, ngày nắng. Dần anh thấy mình giàu. Biết ơn những gì nhỏ — sẽ thấy mình giàu.}
\textit{(He practiced gratitude for a meal, for family, for a sunny day. Gradually he felt rich. Be grateful for the small — you'll feel rich.)}
\end{truyen}

\begin{baihoc}
Biết ơn là giàu.
\textit{(Gratitude is wealth.)}
\end{baihoc}

\begin{ghinhoanh}
\textbf{Short English to remember:}

Gratitude is wealth.

\end{ghinhoanh}

\begin{tuhocnhanh}
1. Câu chuyện này gợi cho bạn suy nghĩ gì? \textit{What does this story make you think?}

2. Bạn sẽ nhớ bài học nào nhất? \textit{Which lesson will you remember most?}
\end{tuhocnhanh}
\ngancach


\section{Đừng vì một chapter xấu mà đóng cả cuốn sách đời mình}
\begin{truyen}{Đừng vì một chapter xấu mà đóng cả cuốn sách đời m...}{Life lesson}
\chuhoa{M}{ột đoạn đời tệ không có nghĩa cả đời tệ. Đừng vì một chapter xấu mà đóng cả cuốn sách đời mình.}
\textit{(One bad stretch doesn't mean the whole life is bad. Don't close the whole book of your life because of one bad chapter.)}
\end{truyen}

\begin{baihoc}
Còn nhiều trang phía trước.
\textit{(Many pages ahead.)}
\end{baihoc}

\begin{ghinhoanh}
\textbf{Short English to remember:}

Many pages ahead.

\end{ghinhoanh}

\begin{tuhocnhanh}
1. Câu chuyện này gợi cho bạn suy nghĩ gì? \textit{What does this story make you think?}

2. Bạn sẽ nhớ bài học nào nhất? \textit{Which lesson will you remember most?}
\end{tuhocnhanh}
\ngancach


\section{Người thật sự yêu thương mình không làm mình nhỏ bé}
\begin{truyen}{Người thật sự yêu thương mình không làm mình nhỏ b...}{Life lesson}
\chuhoa{T}{ình yêu đích thực nâng ta lên, không đè ta xuống. Người thật sự yêu thương mình không làm mình nhỏ bé.}
\textit{(Real love lifts you up, doesn't push you down. Those who truly love you don't make you small.)}
\end{truyen}

\begin{baihoc}
Yêu là nâng đỡ.
\textit{(Love is support.)}
\end{baihoc}

\begin{ghinhoanh}
\textbf{Short English to remember:}

Love is support.

\end{ghinhoanh}

\begin{tuhocnhanh}
1. Câu chuyện này gợi cho bạn suy nghĩ gì? \textit{What does this story make you think?}

2. Bạn sẽ nhớ bài học nào nhất? \textit{Which lesson will you remember most?}
\end{tuhocnhanh}
\ngancach


\section{Sự thật đau nhưng nói dối đau hơn về lâu dài}
\begin{truyen}{Sự thật đau nhưng nói dối đau hơn về lâu dài}{Life lesson}
\chuhoa{A}{nh từng nói dối để tránh đau. Về sau sự thật lộ ra, đau gấp bội. Sự thật đau nhưng nói dối đau hơn về lâu dài.}
\textit{(He once lied to avoid pain. Later when the truth came out, the pain multiplied. Truth hurts but lies hurt more in the long run.)}
\end{truyen}

\begin{baihoc}
Thật thà là an toàn.
\textit{(Honesty is safety.)}
\end{baihoc}

\begin{ghinhoanh}
\textbf{Short English to remember:}

Honesty is safety.

\end{ghinhoanh}

\begin{tuhocnhanh}
1. Câu chuyện này gợi cho bạn suy nghĩ gì? \textit{What does this story make you think?}

2. Bạn sẽ nhớ bài học nào nhất? \textit{Which lesson will you remember most?}
\end{tuhocnhanh}
\ngancach


\section{Mỗi ngày là một cơ hội bắt đầu lại — dù chỉ một chút}
\begin{truyen}{Mỗi ngày là một cơ hội bắt đầu lại — dù chỉ một ch...}{Life lesson}
\chuhoa{A}{nh không cần chờ năm mới. Mỗi ngày anh có thể bắt đầu lại một chút. Mỗi ngày là một cơ hội bắt đầu lại — dù chỉ một chút.}
\textit{(You don't need to wait for the new year. Each day you can start again a little. Every day is a chance to start again — even just a bit.)}
\end{truyen}

\begin{baihoc}
Bắt đầu từ hôm nay.
\textit{(Start from today.)}
\end{baihoc}

\begin{ghinhoanh}
\textbf{Short English to remember:}

Start from today.

\end{ghinhoanh}

\begin{tuhocnhanh}
1. Câu chuyện này gợi cho bạn suy nghĩ gì? \textit{What does this story make you think?}

2. Bạn sẽ nhớ bài học nào nhất? \textit{Which lesson will you remember most?}
\end{tuhocnhanh}
\ngancach


\section{Trưởng thành là khi mình không còn đổ lỗi cho ai}
\begin{truyen}{Trưởng thành là khi mình không còn đổ lỗi cho ai}{Life lesson}
\chuhoa{A}{nh từng đổ lỗi cho bố mẹ, hoàn cảnh, số phận. Trưởng thành là khi mình không còn đổ lỗi cho ai.}
\textit{(He used to blame his parents, circumstances, fate. Growing up is when you stop blaming anyone.)}
\end{truyen}

\begin{baihoc}
Trách nhiệm là trưởng thành.
\textit{(Responsibility is maturity.)}
\end{baihoc}

\begin{ghinhoanh}
\textbf{Short English to remember:}

Responsibility is maturity.

\end{ghinhoanh}

\begin{tuhocnhanh}
1. Câu chuyện này gợi cho bạn suy nghĩ gì? \textit{What does this story make you think?}

2. Bạn sẽ nhớ bài học nào nhất? \textit{Which lesson will you remember most?}
\end{tuhocnhanh}
\ngancach


\section{Bình an không phải không có bão — mà là vững trong bão}
\begin{truyen}{Bình an không phải không có bão — mà là vững trong...}{Life lesson}
\chuhoa{Đ}{ời không phải lúc nào cũng nắng. Bình an là khi bão tới mình vẫn đứng. Bình an không phải không có bão — mà là vững trong bão.}
\textit{(Life isn't always sunny. Peace is when the storm comes and you still stand. Peace isn't no storm — it's standing firm in it.)}
\end{truyen}

\begin{baihoc}
Vững trong sóng gió.
\textit{(Steady in the storm.)}
\end{baihoc}

\begin{ghinhoanh}
\textbf{Short English to remember:}

Steady in the storm.

\end{ghinhoanh}

\begin{tuhocnhanh}
1. Câu chuyện này gợi cho bạn suy nghĩ gì? \textit{What does this story make you think?}

2. Bạn sẽ nhớ bài học nào nhất? \textit{Which lesson will you remember most?}
\end{tuhocnhanh}
\ngancach


\section{Cho người khác cơ hội thứ hai — nhưng đừng quên bài học lần một}
\begin{truyen}{Cho người khác cơ hội thứ hai — nhưng đừng quên bà...}{Life lesson}
\chuhoa{A}{nh có thể tha thứ và cho cơ hội lại. Nhưng anh không quên đã học gì. Cho người khác cơ hội thứ hai — nhưng đừng quên bài học lần một.}
\textit{(You can forgive and give a second chance. But you don't forget what you learned. Give others a second chance — but don't forget the first lesson.)}
\end{truyen}

\begin{baihoc}
Tha nhưng không quên.
\textit{(Forgive but don't forget.)}
\end{baihoc}

\begin{ghinhoanh}
\textbf{Short English to remember:}

Forgive but don't forget.

\end{ghinhoanh}

\begin{tuhocnhanh}
1. Câu chuyện này gợi cho bạn suy nghĩ gì? \textit{What does this story make you think?}

2. Bạn sẽ nhớ bài học nào nhất? \textit{Which lesson will you remember most?}
\end{tuhocnhanh}
\ngancach


\section{Sống đơn giản là biết bỏ bớt — không phải có ít}
\begin{truyen}{Sống đơn giản là biết bỏ bớt — không phải có ít}{Life lesson}
\chuhoa{Đ}{ơn giản không phải nghèo. Đơn giản là biết cái gì đủ. Sống đơn giản là biết bỏ bớt — không phải có ít.}
\textit{(Simple isn't poor. Simple is knowing what's enough. To live simply is to know what to let go — not to have little.)}
\end{truyen}

\begin{baihoc}
Đủ là giàu.
\textit{(Enough is rich.)}
\end{baihoc}

\begin{ghinhoanh}
\textbf{Short English to remember:}

Enough is rich.

\end{ghinhoanh}

\begin{tuhocnhanh}
1. Câu chuyện này gợi cho bạn suy nghĩ gì? \textit{What does this story make you think?}

2. Bạn sẽ nhớ bài học nào nhất? \textit{Which lesson will you remember most?}
\end{tuhocnhanh}
\ngancach


\section{Cuối cùng điều quan trọng không phải ta có gì — mà ta đã cho đi gì}
\begin{truyen}{Cuối cùng điều quan trọng không phải ta có gì — mà...}{Life lesson}
\chuhoa{K}{hi nhìn lại, anh không hỏi mình có bao nhiêu. Anh hỏi mình đã cho đi gì. Cuối cùng điều quan trọng không phải ta có gì — mà ta đã cho đi gì.}
\textit{(When looking back he didn't ask how much he had. He asked what he had given. In the end what matters isn't what we have — but what we gave.)}
\end{truyen}

\begin{baihoc}
Cho đi là còn lại.
\textit{(What we give is what remains.)}
\end{baihoc}

\begin{ghinhoanh}
\textbf{Short English to remember:}

What we give is what remains.

\end{ghinhoanh}

\begin{tuhocnhanh}
1. Câu chuyện này gợi cho bạn suy nghĩ gì? \textit{What does this story make you think?}

2. Bạn sẽ nhớ bài học nào nhất? \textit{Which lesson will you remember most?}
\end{tuhocnhanh}
\ngancach
